\section{Implementation}
\label{sec:implementation}

The composite scheduler is implemented as a request-ordering step that runs
before vLLM's own engine is invoked, rather than as a modification to vLLM
itself. For a batch of prompts, the OPT-125M predictor's score and each
prompt's tokenized length are computed once, combined via
Eq.~\eqref{eq:composite-score}, and used to sort the batch; the sorted
prompts are then submitted to vLLM's offline \texttt{LLM.generate()} API as a
single call, and vLLM's continuous-batching scheduler manages iteration-level
admission and eviction from that point onward. This design keeps the
extension isolated to the scheduling layer and avoids depending on any
internal vLLM APIs that might change between releases.

Two software-environment deviations from the original LTR framework were
necessary to run this pipeline on current hardware and are documented here for
reproducibility. First, the original framework targets vLLM 0.4.1, which has
no prebuilt wheel for the Python 3.12 environment used in this work and does
not build from source in that environment; we instead use a current stable
vLLM release (0.6.1.post2), which also avoids a broken transitive dependency
(\texttt{pyairports}) that affects guided decoding in vLLM 0.5.x. Second, the
serving GPU's pre-installed PyTorch build is not ABI-compatible with vLLM's
compiled CUDA extensions, which required a clean reinstallation of PyTorch as
part of the vLLM install step rather than reusing the platform-provided build.
Both deviations affect only the serving environment and not the scheduling
logic evaluated in this paper, but the underlying vLLM version, being newer
than the original framework's target, is a relevant factor when comparing
absolute latency numbers across studies.

Time-to-first-token (TTFT), reported in Section~\ref{sec:evaluation} as a
secondary indicator of HOL blocking, is computed from vLLM's own per-request
metrics (\texttt{arrival\_time} and \texttt{first\_token\_time}) attached to
each \texttt{RequestOutput} object, rather than approximated from aggregate
batch timing. This metric was added purely additively alongside the existing
throughput and latency instrumentation and does not alter the scheduling
logic or the per-token latency values reported for the four policies.
